\documentclass[10pt,a4paper]{amsart}
\usepackage[margin=19mm]{geometry}
\usepackage{amsmath,amssymb,mathtools}
\usepackage[scr=rsfs,cal=euler]{mathalfa}
\usepackage{xfrac}
\usepackage[unicode,hidelinks,bookmarks=false]{hyperref}
\newcommand{\ie}{i.e.\ }
\newcommand{\cf}{cf.\ }
\newcommand{\resp}{respectively\ }
\newcommand{\Subsection}{\subsection}
\providecommand{\nobreakdash}{\nobreak\hbox{-}}
\setlength{\emergencystretch}{2em}
\multlinegap0pt
\usepackage{bm}
\usepackage{bbm}
\usepackage{dsfont}
\usepackage{xspace}
\usepackage{cancel}
\usepackage{mathtools}
\usepackage{amsthm}
\makeatletter
\@ifundefined{thm}{\newtheorem{thm}{Thm}}{}
\@ifundefined{lem}{\newtheorem{lem}[thm]{Lemma}}{}
\@ifundefined{prop}{\newtheorem{prop}[thm]{Proposition}}{}
\makeatother
\newcommand{\bQ}{\mathbb{Q}}
\newcommand{\bZ} {\mathbb{Z}}
\newcommand{\cU}{\mathcal{U}}
\newcommand{\cX}{\mathcal X}
\newcommand{\cY}{\mathcal Y}
\newcommand{\fm}{\mathfrak m}
\newcommand{\fo}{\mathfrak o}
\newcommand{\lra}{\longrightarrow}
\title[Preperiodic points, finiteness, and structures of semigroups]{Preperiodic points, finiteness, and structures of semigroups of algebraic morphisms}
\author{Jason P. Bell, Thomas J. Tucker}
\date{Source: arXiv:2508.09114v1 (CC BY 4.0)}
\begin{document}
\maketitle

\noindent\emph{Source and licence.} Preperiodic points, finiteness, and structures of semigroups of algebraic morphisms. Version arXiv:2508.09114v1 (CC BY 4.0), distributed under the Creative Commons Attribution 4.0 International licence, as recorded in the arXiv licence metadata for this exact version and shown on the abstract page https://arxiv.org/abs/2508.09114v1. Full rights notice: licenses/arxiv-2508.09114-CC-BY-4.0.txt.\par\smallskip

\noindent\textit{Editorial scope.} Counted unit: the Prop whose source label is \texttt{burn1} in the article (source0.tex lines 446--461, source label \texttt{burn1}), delivered together with the article's own complete original proof of it (source0.tex lines 462--501, one \texttt{proof} environment of 40 lines). No mathematics is written, repaired or re-derived: statement and proof are the article's own text line for line. Required context retained uncounted: basic (lines 359--372). Every definition, display and label of those lines is retained unchanged. Omitted from this package and not redistributed: 1--358, 373--445, 502--1285. Nothing inside a retained excerpt is omitted or altered: every retained body is delivered exactly as the article's lines are written, so no omission receipt of any kind accompanies this package. Citations: the retained excerpts carry no citation command at all, so no bibliography entry of the article is transcribed and no reference is redistributed. Reference adaptation: none was needed and none was made; every reference inside the delivered unit resolves inside it, so no unresolved reference or citation is produced. Printed slips kept literal: no printed word, symbol, hypothesis or number of the counted statement or of its proof was altered. Layout and class: the article's own class is \texttt{amsart} with packages including amsmath, amssymb, comment, xypic; the wrapper is the lane's pinned \texttt{amsart} editorial wrapper at 10pt on A4, which restates the theorem-like environments the retained text needs and adds the article's own macro definitions that the retained text needs (\texttt{\textbackslash bQ}, \texttt{\textbackslash bZ}, \texttt{\textbackslash cU}, \texttt{\textbackslash cX}, \texttt{\textbackslash cY}, \texttt{\textbackslash fm}, \texttt{\textbackslash fo}, \texttt{\textbackslash lra}), with extra wrapper packages (bm, bbm, dsfont, xspace, cancel, mathtools). The delivered numbering is the unit's own. Assets: no retained span contains a figure environment, a graphics-inclusion command, a PDF-page inclusion, a table or a TikZ picture; no image file is copied into this package, the package declares no asset dependency and no image file is redistributed with it. Licence: version arXiv:2508.09114v1 of this article is distributed under the Creative Commons Attribution 4.0 International licence, as recorded in the arXiv licence metadata for this exact version and shown on the abstract page https://arxiv.org/abs/2508.09114v1. A full rights notice is delivered at licenses/arxiv-2508.09114-CC-BY-4.0.txt. Characters: every retained excerpt is pure ASCII, so the delivered engine, XeLaTeX with its default Latin Modern text font, reports no missing character.
\begin{lem}\label{basic}
  Let $\fo_v$ be the integral closure of $\bZ_p$ in some finite extension
  $K_v$ of $\bQ_p$ for some $p$, let $\fm_v$ be the maximal ideal in
  $\fo_v$, and let $k_v$ be the residue field $\fo_v/\fm_v$.  Let
  $\cX$ be a smooth $\bZ_p$-scheme whose generic fiber is a variety
  $X$ over $K_v$.  Let $\gamma \in \cX(k_v)$ and let
  $\cU_\gamma = r^{-1}(\gamma)$, where
  $r: \cX(\fo_v) \lra \cX(k_v)$ is the usual reduction map.   Let
  $G_\gamma$ be a group of automorphisms $\sigma$ of $\cX$ such
  that $\sigma(\cU_\gamma) = \cU_\gamma$.  Then $G_\gamma$ has a
  subgroup $H_\gamma$ of finite index such that $h(W) = W$ for any
  irreducible subvariety $W$ of the generic fiber of $\cX$ such that
  the closure of $W$ in $\cX$ contains a point in $\cU_\gamma$.     
\end{lem}

\begin{prop}\label{burn1}
Let $\fo_v$ be the integral closure of $\bZ_p$ in some finite extension
$K_v$ of $\bQ_p$ for some $p$.  Let $\cX$ be a $\bZ_p$-scheme
whose generic fiber is a variety $X$ over $K_v$, let $G$ be a group
of automorphisms of $\cX$, and let $T$ be the set
of all points in $\cX(\fo_v)$ that are periodic under every element of
$G$.  Let $M$ be the set of all elements of $G$ such that $ht = t$ for
all $t \in T$.  

Then we have the following:
\begin{itemize}
\item [(i)] $M$ is normal in $G$;
\item[(ii)] $[G:M]$ is finite; and
  \item [(iii)] for any $t \in T$, the orbit $Gt$ is finite. 
\end{itemize}
\end{prop}

\begin{proof}
  We see that (iii) follows immediately from (ii) since for any
  $t \in T$, the stabilizer $G_t$ of $t$ contains $M$.  Note that we
  may assume that $X$ is geometrically irreducible in (i) since the
  stabilizer of each component of $X$ has finite index in $G$.
  Likewise, (i) follows easily from (iii) since  if $Gt$ is finite
  for all $t \in T$, then $gt \in T$ for any $g
  \in G$ and $t \in T$, so $ghg^{-1}(t) = g g^{-1} t = t$ for
  any $h \in M$, which means that $g h g^{-1} \in M$ for any $g \in G$
  and $h \in M$.   Hence, it only remains to prove (ii).  
  
We use induction on the dimension of $X$ and Lemma \ref{basic} to prove
(ii).  If $X$
has dimension 0, this is trivial.  Otherwise, we may write $\cX$ as
$\cX' \cup \cY$ where $\cX'$ is smooth over $\bZ_p$ and the generic
fiber of $\cY$ has dimension less than $\dim X$.  An automorphism of a
scheme preserves its smooth locus so $\sigma(\cX') = \cX'$ and
$\sigma(\cY) = \cY$.  We let $M_Y$ be the subgroup of of $G$ that
fixes all elements of $T$ in $\cY$ and $M'$ be the subgroup of $G$
that fixes all elements of $T$ in $\cX'$.  Then $[G:M_Y]$ is finite,
by the inductive hypothesis, so it will suffice to show that $[G:M']$
is finite
since $[G:M_Y \cap M'] \leq [G:M'] \cdot [G:M_Y]$.  

Now, since there are finitely many $\gamma \in \cX'(k_v)$, there are
finitely many residue classes $\cU_\gamma$ of $\fo_v$-points as in
Lemma \ref{basic} containing elements of $T$.  Since $G$ permutes the
residue classes $\cU_\gamma$, there is a subgroup $H$ of finite index
in $G$ such that $h(\cU_\gamma) = \cU_\gamma$ for all $h \in H$ and
all residue classes $\cU_\gamma$. By Lemma \ref{basic}, for each
$\gamma \in \cX'(k_v)$, there is a subgroup $H_\gamma$ of finite index
in $H$ such that $ht =t$ for all $t \in \cU_\gamma$ that are periodic
under all elements of $H$.  Let
$M'' = \cap_{\gamma \in \cX'(k_v)} H_\gamma$.  Now if $t \in T$, then
$t$ is periodic under every element of $H$ and furthermore
$t \in \cU_\gamma$ for some $\gamma \in \cX'(K_v)$; hence, $M''$ fixes
every element of $T$.  Since $M''$ is an intersection of subgroups of
finite index in $G$, it has finite index in $G$.  Clearly, $M'$
contains $M''$ so $[G:M']$ is also finite, as desired.
\end{proof}

\end{document}
